\vspace{\baselineskip}

\centerline{\bf \Huge Олимпиада ЛМШ 2023}
\title{Олимпиада ЛМШ 2023}

\vspace{\baselineskip}
\vspace{\baselineskip}

\problem{1}Начальник решил внедрить Шпиона в организацию НеНачальника. Но чтобы НеРазведывательная служба негодяев не узнала настоящую личность Шпиона, Начальник принял следующие меры безопасности. В архиве, где находится информация о сотрудниках ЛМШ, документы Шпиона положили в специальный инновационный кейс, замок которого выглядит как табличка 6 на 6. Кейс откроется, если во всех клетках таблички будут одинаковые числа. А сами числа изменяются по следующему правилу:\\
1) В начале в правом верхнем углу двойка, а во всех других клетках записаны единички.\\
2)Разрешается за ход прибавить или отнять любое натуральное число либо у любых двух, либо у любых четырех клеток, находящихся на одной диагонали.\\
Смогут ли негодяи узнать личность Шпиона?\\

\problem{2}НеНачальник сумел сбежать от Эльзяты, но во время сражения он выронил 5 монет. По весу все монеты были одинаковыми. Но Начальник ЛМШ знал, что среди них есть хотя бы одна фальшивая монета. Чтобы найти подделки, вызвали Германа с его стендом Ventriloquo, у которого есть способность Embodiment, она позволяет предметам говорить. Но НеНачальник предугадал эту ситуацию и научил свои монеты искусно лгать (настоящие монеты говорят правду, а фальшивые всегда лгут). Монеты положили по кругу и спросили: "Слева от тебя настоящая монета?". Три монеты ответили "Да". Начальник подумал, но не смог определить фальшивки. Тогда Эльзята задала второй вопрос: "Среди твоих соседей есть настоящая монета?". Четыре монеты ответили "Да". И Начальник все понял, но решил проверить Германа. Помогите Герману определить подделки. (Монеты не слепые и естественным образом отличают настоящих от фальшивых) 

\newpage
\vspace{\baselineskip}  
\problem{3}Чтобы раскрыть преступную сеть НеНачальника, Начальник внедрил шпиона НеСашу в качестве одного из помощников НеНачальника. НеСаша обнаружила, что у НеНачальника есть помощники в каждой стране, которые обслуживают секретные туннели, ведущие из одной страны в другую. НеСаша узнала, что в любом тупике находится НеНачальник. (Тупик - место, в которое ведет только один туннель). Кроме того, туннели не образуют циклов, то есть, если выйти из некоторой страны по какому-то туннелю и далее двигаться так, чтобы не проходить по одному туннелю дважды, то невозможно возвратиться в начальную страну. Получив отчет, Начальник заметил несостыковку. И выдвинул гипотезу, что НеНачальников больше одного. Помогите проверить гипотезу.\\

\problem{4}Неначальники научились подделывать монеты всех стран. И чтобы компания ЛМШ не смогла добраться до их главного секрета, они спрятали его в сейфе, закрытый на несколько замков. Каждый НеНачальник носит ключи ко всем замкам, но часть из этих ключей фальшивая. Известно, что среди 10 НеНачальников любые 5 имеют все необходимые истинные ключи. Но никакие 4 НеНачальника не имеют всех настоящих ключей. Сколько же замков охраняют сейф?\\

\problem{5}Открыв Кейс, все вокруг (Эльзята, Начальник, Герман, Расул Владимирович и 5 пойманных НеНачальников) были телепортированы в ловушку. И внутри нее была записка: "Только истинный НеНачальник способен избежать ловушки". Также рядом находилась табличка 100 на 100. Внутри каждой клетки таблицы были фальшивые монеты орлом вверх. Можно одновременно переворачивать монеты во всех клетках одного столбца или строки. Для того чтобы выбраться из ловушки, надо ровно 2006 монет поставить решкой вверх. Можно ли выбраться из этой ловушки?